% !TeX encoding = UTF-8
\documentclass[variant=guia,cover=institutional,mode=screen]{ua-engineering-v6-article}
% !TeX encoding = UTF-8
\UASetUniversity{Universidad Austral}
\UASetFaculty{Facultad de Ingeniería}
\UASetDepartment{Departamento de Computación}
\UASetCampus{Pilar}
\UASetCareer{Ingeniería Informática}
\UASetCommission{Comisión A}
\UASetProfessor{Equipo docente}
\UASetSemester{Segundo cuatrimestre 2026}
\UASetCourse{Sistemas Distribuidos}
\UASetDate{2026}


\UASetClassNumber{01}
\UASetClassTitle{Pensamiento distribuido}
\UASetDocKind{Guía de ejercicios}

\title{Clase 01 --- Pensamiento distribuido}
\author{\UAProfessor}
\date{\UADate}

\begin{document}

\section{Indicaciones generales}
Resuelva cada ejercicio declarando supuestos, modelo de fallas, garantía y trade-off. Cuando corresponda, construya una ejecución verificable y vincule la respuesta con evidencia observable.

\section{Nivel 1: fundamentos y reconocimiento}
\begin{uaexercise}
\textbf{1.}
Redactar una definición operativa de sistema distribuido que no dependa solo de “hay varias máquinas”.
\end{uaexercise}

\begin{uaexercise}
\textbf{2.}
Explicar con tus palabras por qué una demora suficientemente larga puede parecer una falla.
\end{uaexercise}

\begin{uaexercise}
\textbf{3.}
Describir un ejemplo donde el observador remoto no pueda distinguir entre dos historias internas distintas.
\end{uaexercise}

\begin{uaexercise}
\textbf{4.}
Explicar por qué la red no debe modelarse como un canal instantáneo y perfecto.
\end{uaexercise}

\begin{uaexercise}
\textbf{5.}
Indicar cuál es el error conceptual en la afirmación: “si hubo timeout, el servidor se cayó”.
\end{uaexercise}


\section{Nivel 2: ejecuciones y fallas}
\begin{uaexercise}
\textbf{6.}
Construir dos ejecuciones distintas del ejemplo del pago que sean indistinguibles para el cliente.
\end{uaexercise}

\begin{uaexercise}
\textbf{7.}
Modificar el ejemplo del pago para que, además de no llegar la respuesta, exista retry del cliente.
¿Qué nuevo riesgo aparece?
\end{uaexercise}

\begin{uaexercise}
\textbf{8.}
Diseñar un ejemplo de sistema de mensajería donde un mensaje pueda aparecer “no entregado” y, sin embargo, haber sido procesado en el servidor.
\end{uaexercise}

\begin{uaexercise}
\textbf{9.}
Describir un caso donde el problema no sea la caída de un nodo, sino la pérdida de conectividad entre nodos correctos.
\end{uaexercise}

\begin{uaexercise}
\textbf{10.}
Explicar por qué una respuesta tardía puede ser más difícil de manejar que una falla explícita.
\end{uaexercise}


\section{Nivel 3: diseño y comparación}
\begin{uaexercise}
\textbf{11.}
Para el caso del pago con timeout, completar una tabla con tres columnas: síntomas relevados, causas posibles y evidencias necesarias para distinguirlas.
\end{uaexercise}

\begin{uaexercise}
\textbf{12.}
Separar el análisis del caso anterior en tres partes: qué se sabe, qué se infiere y qué todavía no puede afirmarse.
\end{uaexercise}

\begin{uaexercise}
\textbf{13.}
Proponer una respuesta de diseño para el caso anterior conectando relevamiento, análisis y propuesta.
La respuesta debe indicar garantía buscada, solución elegida y trade-off aceptado.
\end{uaexercise}

\begin{uaexercise}
\textbf{14.}
Aplicar el esquema $D=(P,F,G,S,T)$ a un sistema de compra online.
\end{uaexercise}

\begin{uaexercise}
\textbf{15.}
Aplicar el esquema $D$ a una plataforma de chat.
\end{uaexercise}


\section{Nivel 4: integración y defensa}
\begin{uaexercise}
\textbf{16.}
Aplicar el esquema $D$ a un sistema de autenticación distribuido.
\end{uaexercise}

\begin{uaexercise}
\textbf{17.}
Comparar dos soluciones distintas para el mismo problema y mostrar cómo cambia el trade-off.
\end{uaexercise}

\begin{uaexercise}
\textbf{18.}
Diseñar una estrategia conceptual para evitar pagos duplicados cuando el cliente reintenta automáticamente.
\end{uaexercise}

\begin{uaexercise}
\textbf{19.}
Explicar qué costo introduce volver idempotente una operación que naturalmente no lo es.
\end{uaexercise}

\begin{uaexercise}
\textbf{20.}
Proponer una garantía razonable para una red social bajo partición y justificarla.
\end{uaexercise}


\end{document}
